\documentclass[conference]{IEEEtran}
\usepackage[utf8]{inputenc}
\usepackage{hyperref}
\title{What are the effects of intermittent fasting on cognitive performance?}
\author{\IEEEauthorblockN{ResearchGenie}\IEEEauthorblockA{AI-generated research draft}}
\begin{document}
\maketitle
\textit{\textbf{AI-generated research draft for human review.} Not a verified or published research paper. Every claim must be checked against its cited source before any external use.}\par

\section{Introduction}
Intermittent fasting (IF) is an eating pattern defined by periodic cycles of fasting and feeding that promotes health and cognitive improvement [1]. This approach functions partly through the microbiota-gut-brain axis, a communication bridge between gut microbiota and the brain that influences cognitive function via the immune, endocrine, and nervous systems [1]. However, visceral adipose tissue acts as an active endocrine organ secreting proinflammatory cytokines such as TNF-α, IL-6, and CRP, which contribute to systemic low-grade inflammation, a hallmark of obesity-related metabolic dysfunction [2]. In a randomised, waitlist-controlled trial involving 122 participants adapting to a 16-hour breakfast-skipping intermittent fasting protocol over ten days, Bayesian multi-level modelling revealed that fasting participants did not exhibit lower cognitive performance or mood compared to control participants [3].

The literature review synthesizes evidence on intermittent fasting (IF) and cognitive performance across human trials, animal models, and mechanistic reviews. Human data suggests that IF effects are driven by negative energy balance and weight loss rather than the fasting protocol itself, while systemic inflammation (e.g., IL-6) is a critical determinant of cognitive decline [2]. Animal studies demonstrate neuroprotective mechanisms involving reduced neuroinflammation and enhanced synaptic plasticity [4][5]. Adaptation studies show stable objective performance but subjective concentration deficits, whereas Ramadan fasting contexts reveal progressive declines in young athletes linked to sleep disruption [6][7]. Mechanistic reviews highlight the role of the gut-brain axis, noting potential synergies between IF and probiotics [1].

The primary analytical challenge in evaluating intermittent fasting (IF) and cognitive performance is distinguishing whether observed effects stem from the temporal pattern of eating or from associated metabolic changes like weight loss. Evidence suggests that when energy balance is controlled, the specific fasting-feeding cycle may not be the sole driver of cognitive outcomes [2]. Instead, systemic inflammation driven by adipose tissue appears to influence cognitive decline by compromising the blood–brain barrier, increasing permeability to cytokines and neurotoxins in regions such as the hippocampus [2]. Consequently, metabolic shifts resulting from body fat loss—specifically reductions in inflammation and insulin resistance—likely contribute to cognitive improvements observed with IF and caloric restriction [2]. However, these relationships describe associations rather than causation; for instance, improved insulin resistance may be a concurrent benefit of fat loss rather than a direct mediator of cognitive function [2].

Methodological heterogeneity across studies constrains the ability to draw definitive conclusions regarding intermittent fasting (IF) and cognitive performance. A primary barrier is the lack of standardization in outcome measures; included reviews have grouped diverse cognitive tests and domains into single categories, preventing comparisons between IF protocols or even within similar interventions [2]. This aggregation obscures potential variations in susceptibility across specific cognitive functions. Furthermore, analytical approaches often rely on categorical yes/no comparisons rather than correlating the magnitude of physiological changes—such as body weight or fat loss—with performance metrics [2]. Such a binary approach precludes understanding whether energy restriction or adiposity loss specifically accounts for observed effects. Additionally, distinct IF modalities, including alternate-day fasting, time-restricted feeding, and 5:2 protocols, were frequently collapsed into a single category, inhibiting analysis of how specific forms influence cognition [2]. Temporal resolution in assessment also introduces uncertainty regarding the timing of cognitive shifts, as objective measures often occur at the end of the fasting window, potentially missing within-day changes [3].

\section{Background}
Intermittent fasting (IF) is defined as an eating pattern that promotes health and cognitive improvement through periodic cycles of fasting and feeding [1]. This approach functions by modulating the microbiota-gut-brain axis (MGBA), a communication bridge between gut microbiota and the brain that influences cognitive function via the immune, endocrine, and nervous systems [1]. Visceral adipose tissue acts as an active endocrine organ secreting proinflammatory cytokines such as TNF-α, IL-6, and CRP, which contribute to systemic low-grade inflammation, a hallmark of obesity-related metabolic dysfunction [2]. In a randomised, waitlist-controlled trial involving 122 participants adapting to a 16-hour breakfast-skipping intermittent fasting protocol over ten days, Bayesian multi-level modelling revealed that fasting participants did not exhibit lower cognitive performance or mood compared to control participants [3].

\section{Literature Review}
Human trial evidence indicates that reductions in visceral adiposity and improvements in insulin resistance are key factors linked to enhanced cognitive performance [2]. Comparative analyses of intermittent fasting (IF) trials against caloric restriction (CR) trials reveal comparable cognitive outcomes, supporting the inference that negative energy balance, rather than specific fasting protocols or altered eating patterns, drives these effects [2]. Achievement of adequate weight or body fat loss in adults with overweight or obesity, whether through IF or CR, is associated with changes in cognitive performance across clinical trials [2].

Beyond metabolic mechanisms, systemic inflammation emerges as a critical determinant of cognitive trajectory. Elevated IL-6 levels were significantly associated with slower processing speed and poorer executive function in the Sydney Memory and Ageing Study, even after controlling for depression, vascular risk factors, and APOE-ε4 status [2]. Similarly, subjects in the Health, Aging, and Body Composition Study who exhibited high tertiles of IL-6 and C-reactive protein performed significantly worse on cognitive tests at follow-up and faced higher odds of cognitive decline over two years [2]. In older Dutch adults without dementia, higher IL-6 levels were associated with greater cognitive decline across domains such as memory and processing speed over a 3-year period [2]. Despite differences in sample demographics and geographic regions, all three cited studies reported worse cognitive outcomes in groups with the highest inflammatory burden compared to those with the lowest [2].

Animal models further elucidate the neuroprotective potential of IF against specific insults. In 18-month-old male mice subjected to abdominal exploration under sevoflurane anesthesia, a preoperative intermittent fasting regimen of 16 hours fasting and 8 hours eating for two weeks significantly enhanced learning and memory capabilities compared to control groups [4]. This enhancement was measured using novel object recognition and Morris water maze tests, showing marked improvements following the intervention [4]. Mechanistically, IF was associated with increased synaptic plasticity, evidenced by improved function in hippocampal neurons of aged mice following anesthesia [4]. The administration of intermittent fasting markedly enhanced the expression of synaptic-associated proteins in hippocampal neurons, concomitant with decreased expression of pro-inflammatory factors and reduced microglial cell density within the hippocampus [4]. In a rat model of systemic bacterial infection induced by intravenous lipopolysaccharide (LPS), a 30-day intermittent fasting regimen ameliorated cognitive deficits in the Barnes maze and inhibitory avoidance tests compared to rats on an ad libitum diet, without altering general locomotor activity or motor performance [5]. The cognitive protection observed was associated with reduced hippocampal expression of the LPS receptor TLR4 and inducible nitric oxide synthase (iNOS) following LPS administration [5]. IF also prevented the LPS-induced elevation of pro-inflammatory cytokines (IL-1α, IL-1β, TNF-α) in both the hippocampus and serum, while simultaneously increasing levels of the anti-inflammatory cytokine IL-10 in the hippocampus [5]. Furthermore, in the setting of systemic inflammation, intermittent fasting prevented the LPS-induced reduction of Brain-Derived Neurotrophic Factor (BDNF) levels in the hippocampus, thereby sustaining neurotrophic support critical for synaptic plasticity and cognitive function [5].

Human studies examining adaptation to IF protocols report stable objective cognitive performance during the initial phase. In a randomised, waitlist-controlled trial involving N=122 participants adapting to a 16-hour breakfast-skipping intermittent fasting protocol over ten days, Bayesian multi-level modelling revealed that fasting participants did not exhibit lower cognitive performance or mood compared to control participants [3]. While objective cognitive performance and mood remained stable during the adaptation period, participants consistently reported feeling less concentrated while fasting before noon compared to after breaking the fast in the afternoon [3]. Positive mood was more strongly positively associated with subjective cognitive performance in the fasting group compared to the control group, suggesting that psychological factors may play a prominent role in linking energy intake and cognitive function beyond metabolic mechanisms [3].

Conversely, specific contexts such as religious fasting during Ramadan present distinct challenges. In young female handball players, baseline morning-afternoon differences in attention, simple reaction time, squat jump, Illinois agility, and total/peak 5-m shuttle run distances were significantly greater in the afternoon compared to the morning during the ten days before Ramadan [6]. During the first ten days of Ramadan, these previously observed morning-afternoon differences in physical and cognitive performances were reduced compared to the pre-Ramadan baseline [6]. However, during the last ten days of Ramadan, morning-afternoon differences in physical and cognitive performances were further blunted or reversed relative to baseline [6]. In adolescent female handball players, Ramadan intermittent fasting was associated with a significant decline in cognitive and psychomotor performance, particularly during the fourth week of fasting [7]. This decline coincided with a gradual decrease in sleep quality, duration, and Pittsburgh Sleep Quality Index (PSQI) scores throughout the Ramadan period, reaching their lowest levels in the fourth week of fasting [7].

Mechanistic reviews suggest that intermittent fasting enhances neuroplasticity and reduces oxidative stress and inflammation [1]. Probiotic supplementation enhances cognitive performance by modulating gut microbiota composition and increasing short-chain fatty acid production [1]. The modulation of gut microbiota composition and the increase in short-chain fatty acid production by probiotics promote neurogenesis and synaptic plasticity [1]. Intermittent fasting enhances gut microbial diversity and metabolic efficiency [1], while probiotics modulate gut barrier integrity and neurotransmitter synthesis [1]. The combination of probiotics and intermittent fasting may exert complementary effects on cognitive function, though this remains speculative without direct empirical evidence [1].

\section{Discussion}
The primary analytical challenge in evaluating intermittent fasting (IF) and cognitive performance is distinguishing whether observed effects stem from the temporal pattern of eating or from associated metabolic changes like weight loss. Evidence suggests that when energy balance is controlled, the specific fasting-feeding cycle may not be the sole driver of cognitive outcomes [2]. Instead, systemic inflammation driven by adipose tissue appears to influence cognitive decline by compromising the blood–brain barrier, increasing permeability to cytokines and neurotoxins in regions such as the hippocampus [2]. Consequently, metabolic shifts resulting from body fat loss—specifically reductions in inflammation and insulin resistance—likely contribute to cognitive improvements observed with IF and caloric restriction [2]. However, these relationships describe associations rather than causation; for instance, improved insulin resistance may be a concurrent benefit of fat loss rather than a direct mediator of cognitive function [2].

Research in specific populations further clarifies the role of adaptation. In young female handball players, Ramadan intermittent fasting did not necessarily impair intraday cognitive performance, suggesting that brief fasting (up to 18 hours) does not inherently harm cognition due to energy availability, as the brain maintains access to resources and utilizes ketones [3]. Conversely, when deterioration in cognitive function was observed during Ramadan, it was primarily attributed to cumulative fatigue and reduced sleep quality rather than dietary intake or body composition changes [7]. These findings imply that IF may be a viable health intervention for cognitively demanding tasks provided an adequate adaptation period is respected [3].

Mechanistic studies in animal models highlight potential neuroprotective pathways. Intermittent fasting has been shown to mitigate inflammation in the hippocampal area, potentially safeguarding against cognitive impairment and synaptic plasticity loss induced by anesthesia [4]. In models of lipopolysaccharide-induced neuroinflammation, IF shifted gene expression toward anti-inflammatory signaling cascades mediated by the preservation of inhibitory proteins like IκB, resulting in lower overall inflammatory responses despite transient increases in NF-κB activation [5]. These systemic effects extended beyond the central nervous system, with IF inhibiting increases in serum levels of TNF-α, IL-6, and other pro-inflammatory cytokines [5].

Despite these mechanistic insights, critical gaps remain regarding the interplay between IF and other interventions. The regulatory mechanisms governing how probiotics and intermittent fasting jointly influence cognitive function via the microbiota–gut–brain axis represent a therapeutic area that remains underexplored in current research [1]. Future directions must therefore focus on evaluating IF interventions in human subjects at risk for neurological conditions involving neuroinflammation, leveraging its demonstrated ability to suppress inflammatory processes and sustain hippocampal BDNF levels [5].

\section{Research Gap}
\subsection{Absence of Human Trials for Dementia Prevention}
The corpus contains no human trials investigating IF specifically for the prevention or mitigation of dementia in at-risk populations. Paper [1] explicitly notes that while IF has been studied in hypertension and diabetes, it has not yet been explored as a strategy to reduce the risk of emergent dementia in humans. Papers [4] and [8] provide mechanistic evidence in rats and mice regarding neuroinflammation and postoperative dysfunction but do not translate these findings to human clinical populations with cognitive decline.

\textit{Evidence:} Paper [1] states: 'In humans, intermittent fasting has been studied in hypertension, diabetes, and related conditions, but, to date, not as a strategy to reduce the risk of emergent dementia.' Papers [4] and [8] utilize rat and mouse models respectively, leaving a translational gap for human aging populations.

\subsection{Lack of Within-Day Resolution in Human Cognitive Assessments}
Current human studies rely on between-subject or between-day comparisons, missing acute cognitive dynamics. Paper [6] notes that while fasting participants felt less concentrated before noon, objective performance did not differ from controls, suggesting a potential dissociation between subjective and objective measures that is masked by low temporal resolution. Papers [3] and [7] assess athletes at fixed morning and afternoon time points but do not capture the trajectory of cognitive decline or recovery within a single fasting day.

\textit{Evidence:} Paper [6] reports: 'Participants felt less concentrated while fasting, before noon, compared to after breaking the fast... Bayesian multi-level modelling revealed that fasting participants did not have lower cognitive performance.' Paper [7] compares morning vs. afternoon but lacks continuous monitoring.

\subsection{Conflation of IF Mechanisms with Caloric Restriction Effects}
The literature conflates the effects of intermittent fasting (IF) with general caloric restriction (CR). Paper [5] argues that cognitive benefits may be mediated by weight loss and insulin sensitivity rather than the specific fasting pattern, yet most studies do not control for this confound. Without isolating the 'fasting effect' from the 'energy deficit effect,' it is unclear if IF offers unique neuroprotective mechanisms beyond simple calorie counting.

\textit{Evidence:} Paper [5] states: 'Evidence from human trials suggests that reductions in body weight... and improvements in insulin resistance are key factors linked to enhanced cognitive performance.' It further notes that 'further investigation of specific types of IF patterns... may identify mechanisms through which IF can yield benefits on neurocognitive health beyond those of loss of body fat.'

\subsection{Scarcity of Data on Cognitive Performance During Physical Exertion}
Studies involving athletes during Ramadan fasting assess cognitive function only at rest. Paper [3] explicitly highlights this limitation, noting that assessing performance 'under resting conditions rather than during physical exertion' prevents understanding how metabolic stress from fasting interacts with the physiological demands of exercise on cognition.

\textit{Evidence:} Paper [3] states: 'Furthermore, all assessments were conducted under resting conditions rather than during physical exertion. In this aspect, it would be valuable to explore how RIF affects the cognitive function measurements of athletes while they are actively exercising.'

\section{Proposed Novelty and Contribution}
Research addressing these gaps would be novel because it directly targets the methodological and population limitations identified in the corpus. Specifically, studies that isolate IF effects from caloric restriction (Gap 3) or utilize within-day high-resolution monitoring (Gap 2) would provide new mechanistic insights not currently available.

\textit{Caveats:} Novelty is assessed relative to the provided corpus. It is possible that recent unpublished preprints or non-English studies exist that address these questions but were not captured in the retrieval.

\section{Limitations}
Methodological heterogeneity across studies constrains the ability to draw definitive conclusions regarding intermittent fasting (IF) and cognitive performance. A primary barrier is the lack of standardization in outcome measures; included reviews have grouped diverse cognitive tests and domains into single categories, preventing comparisons between IF protocols or even within similar interventions [2]. This aggregation obscures potential variations in susceptibility across specific cognitive functions. Furthermore, analytical approaches often rely on categorical yes/no comparisons rather than correlating the magnitude of physiological changes—such as body weight or fat loss—with performance metrics [2]. Such a binary approach precludes understanding whether energy restriction or adiposity loss specifically accounts for observed effects. Additionally, distinct IF modalities, including alternate-day fasting, time-restricted feeding, and 5:2 protocols, were frequently collapsed into a single category, inhibiting analysis of how specific forms influence cognition [2].

Temporal resolution in assessment also introduces uncertainty regarding the timing of cognitive shifts. In randomized controlled trials, objective performance measures were typically timed to coincide with the end of the fasting window, potentially missing transient within-day fluctuations that occur before meal re-introduction [3]. This design choice limits the capacity to characterize dynamic cognitive trajectories during the fast. Adherence verification presents further challenges; without continuous blood-glucose monitoring or other objective tools, studies cannot fully confirm protocol compliance or distinguish between intervention effects and non-adherence in control groups due to participant burden concerns [3]. Finally, some investigations lack control groups entirely or fail to account for confounding variables such as pre-dawn meal timing and fluid intake, which restricts causal inference regarding IF effects on psychomotor and cognitive skills in specific populations like adolescent athletes [7].

\section{Conclusion}
The current evidence indicates that intermittent fasting (IF) does not inherently impair cognitive performance during adaptation phases. In controlled trials where energy balance is maintained, participants exhibit stable objective performance and mood comparable to controls, suggesting the fasting protocol itself is not a primary driver of cognitive decline [3]. However, observed improvements in cognition are likely secondary to metabolic shifts rather than the temporal eating pattern alone. Systemic inflammation driven by visceral adipose tissue compromises the blood–brain barrier and promotes neurotoxicity; therefore, IF-associated weight loss and reductions in inflammatory markers such as IL-6 may mediate cognitive benefits [2]. Consequently, when energy restriction is absent or controlled, the specific fasting cycle appears insufficient to generate robust cognitive gains independent of metabolic change.

Despite these mechanistic insights, definitive conclusions remain constrained by significant methodological heterogeneity. The lack of standardized outcome measures and the aggregation of diverse cognitive domains into single categories obscure variations in susceptibility across different functions [2]. Furthermore, analytical approaches frequently rely on binary comparisons rather than correlating physiological changes with performance metrics, precluding a clear understanding of whether energy restriction or adiposity loss specifically accounts for observed effects [2]. Distinct IF modalities are often collapsed into single categories, inhibiting analysis of how specific protocols influence cognition, while temporal resolution in assessment may miss within-day cognitive shifts that occur during fasting windows [3].

Future research must prioritize standardizing outcome measures and correlating physiological magnitude with performance to disentangle the effects of energy balance from the fasting protocol itself. Studies should explicitly model inflammatory markers and adiposity changes as mediators to determine if they causally drive cognitive improvements or merely correlate with them. Additionally, investigations need to assess specific IF modalities separately and evaluate cognitive performance at multiple time points within the daily cycle to capture transient effects that current end-point assessments miss.

\begin{thebibliography}{99}
\bibitem{ref1} [1] Y. Chen, C. T. Ho, and X. Zhang, "The regulatory mechanism of intermittent fasting and probiotics on cognitive function by the microbiota‐gut‐brain axis," \textit{Journal of Food Science}, 2025, doi: 10.1111/1750-3841.70132.
\bibitem{ref2} [2] F. Draicchio and K. V. Axen, "Does Energy Restriction and Loss of Body Fat Account for the Effect of Intermittent Fasting on Cognitive Function?," \textit{Nutrients}, 2025, doi: 10.3390/nu17152407.
\bibitem{ref3} [3] C. Bamberg, J. Reichenberger, and J. B. A, "Stable cognitive performance while adapting to intermittent fasting: A randomised controlled trial," \textit{Journal of Health Psychology}, 2025, doi: 10.1177/13591053251351204.
\bibitem{ref4} [4] L. Wang, Q. Wang, and X. Wang, "Intermittent fasting alleviates postoperative cognitive dysfunction by reducing neuroinflammation in aged mice," \textit{Brain Research Bulletin}, 2024, doi: 10.1016/j.brainresbull.2024.111034.
\bibitem{ref5} [5] A. R. Vasconcelos, L. Yshii, and T. A. Viel, "Intermittent fasting attenuates lipopolysaccharide-induced neuroinflammation and memory impairment," \textit{Journal of Neuroinflammation}, 2014, doi: 10.1186/1742-2094-11-85.
\bibitem{ref6} [6] H. Bougrine, M. Cherif, and H. Chtourou, "Does Ramadan intermittent fasting affect the intraday variations of cognitive and high-intensity short-term maximal performances in young female handball players?," \textit{Biological Rhythm Research}, 2023, doi: 10.1080/09291016.2023.2198794.
\bibitem{ref7} [7] H. Bougrine, N. Chalghaf, and C. Azaiez, "The impact of intermittent fasting during Ramadan on psychomotor and cognitive skills in adolescent athletes," \textit{Frontiers in Sports and Active Living}, 2024, doi: 10.3389/fspor.2024.1362066.
\end{thebibliography}

\end{document}
